  \hypertarget{index_intro_sec}{}\section{Introduction}\label{index_intro_sec}
Welcome to a brief introduction of a project executed by a student of the university of Kassel. ~\newline
 The goal of this project is to create reliable software that controlls the robot to drive three laps on a track (\href{track_img.jpg}{\tt {\bfseries example track}}) consisting of a black line, autonomously. ~\newline
There is a black square placed on the track which acts as the startfield. ~\newline
 The robot has three sensors which enable it to detect the black line aswell as the startfield so it can take action accordingly. ~\newline
~\newline
In addition to that the robot starts driving only after being placed on the startfield which initiates a 15s countdown. ~\newline
Upon being moved the countdown shall be interrupted and the robot has to look for the startfield again. ~\newline
 The countdown has the robot blinking three leds with 5hz which are placed on top of the roboter. ~\newline
When the robots completes 3 laps it stops and a 60s countdown starts, analog to the first countdown.\hypertarget{index_install_sec}{}\section{Installation}\label{index_install_sec}
How to compile and install the firmware on the roboter\+:
\begin{DoxyItemize}
\item Before the firmware can be flashed, the robot has to be \href{flashing_process.jpg}{\tt {\bfseries connected}} to a pc containing the firmware.
\item The of process compiling and flashing is automated using Makefile. ~\newline
 To install the default configuration it is sufficient to navigate to the projects directory in the console and run {\bfseries make}
\end{DoxyItemize}\hypertarget{index_Explanation}{}\subsection{Explanation}\label{index_Explanation}

\begin{DoxyItemize}
\item In order to install new firmware on the robot. We first need to \char`\"{}translate\char`\"{} the sourcecode for the computer/ hardware to understand. ~\newline
 This is done by compiling the sourcecode into bytecode. Since the project consists of multiple sourcefiles these also have to be linked together so the output is one object file. ~\newline

\item The project can be compilated in different ways for more information view section {\bfseries Compilation Modes}. ~\newline

\item The {\bfseries A\+T\+M\+E\+G\+A328P} uses a different encoding so the object file has to converted to a hex file ~\newline

\item Now the correct hex file is on the computer and only has to be installed on the microcontroller.
\item This is done by flashing.
\item The hex and flash commands are listed in the {\bfseries Makefile}.
\end{DoxyItemize}\hypertarget{index_comp}{}\subsection{Compilation Modes}\label{index_comp}
The Compilation modes can be adjusted as wanted by setting the values of the {\bfseries -\/ D Flags}. ~\newline
 in the {\bfseries Makefile} under {\bfseries C\+F\+L\+A\+GS}. If a value is set to 0 the corrosponding mode is disabled otherwise it is enabled. ~\newline

\begin{DoxyItemize}
\item {\bfseries D\+E\+B\+U\+G\+\_\+\+M\+O\+DE}\+: ~\newline
 Turns on debug prints that are send to the {\bfseries serial port} via \hyperlink{usart_8c_a2744fce10359c4e63ba8fcb65e7d92d2}{transmit\+\_\+debug\+\_\+msg()} ~\newline
 The data send contains the current State, the sensor values and speed of the engines.
\item {\bfseries C\+O\+U\+N\+T\+D\+O\+N\+W\+\_\+\+M\+O\+DE}\+: ~\newline
 Enables the 15s Countdown at the start. ~\newline
 Otherwise the Roboter starts with the linefollowing algorithm as soon as it is turned on.
\item {\bfseries L\+A\+P\+C\+O\+U\+N\+T\+E\+R\+\_\+\+M\+O\+DE}\+: ~\newline
 If enabled the robot counts the laps otherwise it follows the line till it\textquotesingle{}s batteries die.
\item {\bfseries L\+A\+PS}\+: ~\newline
 Determines how many laps the robot has to drive. ~\newline
 {\bfseries Note}\+: This is only relevant if {\bfseries L\+A\+P\+C\+O\+U\+N\+T\+E\+R\+\_\+\+M\+O\+DE} is enabled.
\end{DoxyItemize}\hypertarget{index_Hardware}{}\section{Hardware}\label{index_Hardware}
Microcontroller Arduino Modell\+: {\bfseries A\+T\+M\+E\+G\+A328P} ~\newline
The \href{pin_layout.pdf}{\tt {\bfseries Pin layout}} can be looked up here.\hypertarget{index_Features}{}\section{Features}\label{index_Features}
The hole implementation is based on the api \hyperlink{fsm_8h}{fsm.\+h} which is a finite state machine. ~\newline
This has several advantages. The robot knows at all times in which state it currently is in. ~\newline
This allows for more accurate state transitions which results in better driving behaviour. ~\newline
Each state can have its own logic and transitions which can be completly independant of each other. ~\newline
This makes it incredibly easy to debug and implement new states. ~\newline
 When adding, a new state the logic of the allready existing transitions does not have to be modified because only the information relevant to the current state has to be interpreted.\hypertarget{index_fsm_explanation}{}\subsection{F\+S\+M Explanation}\label{index_fsm_explanation}

\begin{DoxyItemize}
\item Each State implements an {\bfseries enter function} and an {\bfseries update function} which defines the behaviour of the State. ~\newline

\item The {\bfseries enter function} is called once for every transition into an other state which invokes the new current state\textquotesingle{}s {\bfseries enter function} implementation. ~\newline
 E.\+g. When the State Forward is entered the motor settings can be adjusted accordingly in the \hyperlink{states_8c_a639a5239b6bf87b7629b5a5a781f0623}{enter\+\_\+forward()} function as this only need to happen once in this state.
\item The {\bfseries update function} is called repeadetly in \hyperlink{fsm_8c_a053b1146a3f1c34576d89d599e34329a}{start\+\_\+fsm\+\_\+cycle()} which invokes the current\+\_\+state {\bfseries update function}. ~\newline
 Each State uses this function to do task that have to take place over a period of time and to determine wether to stay or to transition into another State if a certain condition is met.
\item To transition to another state \hyperlink{fsm_8c_a80016823d020a335a595bcd09bc39c96}{transition\+\_\+to\+\_\+state()} is called. This is typically done inside of a States\textquotesingle{}s {\bfseries update function} as it has the needed parameters.
\item {\bfseries Note}\+: Not every State has to implement an {\bfseries enter function} but each State has to implement an {\bfseries update function}.
\end{DoxyItemize}\hypertarget{index_implement_new_states}{}\subsection{How to implement new States}\label{index_implement_new_states}
Descrition of how to add a new state.
\begin{DoxyItemize}
\item Add a new state in State
\item Create function enter\+\_\+ + \char`\"{}name\+\_\+of\+\_\+state\char`\"{} and update\+\_\+ + \char`\"{}name\+\_\+of\+\_\+state\char`\"{} for better readabilty.
\item Call to \hyperlink{fsm_8c_a2d99213b7ce90d5e4b0cbd7f261d8e3e}{add\+\_\+state()} passing a pointer to the States\textquotesingle{}s enter and update function
\item Implement the enter and update function as wanted, achieving the desired behaviour for that State.
\end{DoxyItemize}\hypertarget{index_additional_features}{}\section{Additional Features}\label{index_additional_features}
The {\bfseries A\+T\+M\+E\+G\+A328p} comes with 3 onboard timers which are all used in order to prevent software delays. ~\newline
All delays including the delay for the start and end countdown are implemented via the usage of interrupts. ~\newline
This is more efficient as the microcontroller does not have to waist energy doing nothing. Another advantage is that the programm can still be running normally during the delays which enables us to take measurements as per usuall and take action accordingly.\hypertarget{index_States}{}\section{States}\label{index_States}
The State transitions are implemented as shown in \href{transition_map.pdf}{\tt {\bfseries Transition Map}}.
\begin{DoxyItemize}
\item {\bfseries I\+N\+IT}\+: ~\newline
 Inilizes all the modules.
\item {\bfseries C\+H\+E\+C\+K\+\_\+\+S\+T\+A\+R\+T\+P\+OS}\+: ~\newline
 Lights leds based on their status and transition to C\+H\+E\+C\+K\+\_\+\+S\+T\+A\+R\+T\+P\+OS when the startfield is detected.
\item {\bfseries C\+O\+U\+N\+T\+D\+O\+WN}\+: ~\newline
 Starts the 15s countdown. If the robot is moved during this time the countdown is interrupted. ~\newline
 If the countdown finished successfully the roboter transitions to L\+E\+A\+V\+E\+\_\+\+S\+T\+A\+RT
\item {\bfseries L\+E\+A\+V\+E\+\_\+\+S\+T\+A\+RT}\+: ~\newline
 In this state the robot starts to drive forward until it leaves the startfield. ~\newline
 This State is used so the Robot doesn\textquotesingle{}t count the startfield multiple times per pass.
\item {\bfseries F\+O\+R\+W\+A\+RD}\+: ~\newline
 The Robot drives forward by setting both engines forward and to the same speed. ~\newline
 This state is left only if one or both of the side sensors detect a line.
\item {\bfseries L\+E\+F\+T\+\_\+\+S\+O\+FT}\+: ~\newline
 In this state the robot performes a slight turn by making the left eng speed slower then the right via \hyperlink{motor__controller_8c_aedcc8ddfc6f7a9d3f55833ac09fa01f5}{set\+\_\+direction()} ~\newline
 This state is used mainly to react to slight inaccuracies when the robot wants to drive forward but ~\newline
 is not doing this perfectly. There is always a small margin. This keep the robot on track.
\item {\bfseries R\+I\+G\+H\+T\+\_\+\+S\+O\+FT}\+: ~\newline
 Analog to L\+E\+F\+T\+\_\+\+S\+O\+FT
\item {\bfseries L\+E\+F\+T\+\_\+\+H\+A\+RD}\+: ~\newline
 This state is used for the sharp turnes. By setting the left eng backwards and the right eng to a high speed ~\newline
 the robot performs a sharp turn.
\item {\bfseries R\+I\+G\+H\+T\+\_\+\+H\+A\+RD}\+: ~\newline
 Analog to L\+E\+F\+T\+\_\+\+H\+A\+RD
\item {\bfseries C\+H\+E\+C\+K\+\_\+\+L\+AP}\+: ~\newline
 State that is entered when all 3 sensors see black. ~\newline
 In this case the robot checks if it is indeed on the startfield by checking if all 3 sensors detect black for an extended amount of time. /n It is not sufficient to assume it is on the startfield automatically as it can occure in the sharp turns ~\newline
 that all 3 sensors see black for a very short amount of time. ~\newline
 This state is used to differentiate between these two cases.
\item {\bfseries G\+O\+A\+L\+\_\+\+R\+E\+A\+C\+H\+ED}\+: ~\newline
 The goal is reached when the robot completed all laps. A 60s countdown is initiated.
\end{DoxyItemize}\hypertarget{index_Calibration}{}\section{Calibration}\label{index_Calibration}
Before letting the robot drive it can be wise to calibrate it\textquotesingle{}s motors and sensors. ~\newline
Sometimes one motor can be stronger than the other. So the robot is thinking its driving straight ~\newline
but in reality driving to the site. The same can be said about the sensors. ~\newline
The sensors are not always 100\% accurate so it can be reasonable to adjust the enum Threshold enum as needed. ~\newline
 Each of the three sensors has its own value. ~\newline
\href{https://youtu.be/au92yQbTNe0}{\tt {\bfseries Here}} is an example how the robot should drive when it is calibrated properly.\hypertarget{index_further_information}{}\section{Additional information}\label{index_further_information}
If you wish to read more about the code and why certain code is used the way it is, just explore the option given above in the main page. ~\newline
If you have any suggestions for improvement, do not hesitate and contact me\+: \href{mailto:findrunas@hotmail.com}{\tt findrunas@hotmail.\+com} ~\newline
